\documentclass[10pt,a4paper]{amsart}
\usepackage[margin=19mm]{geometry}
\usepackage{amsmath,amssymb,mathtools}
\usepackage[scr=rsfs,cal=euler]{mathalfa}
\usepackage{xfrac}
\usepackage[unicode,hidelinks,bookmarks=false]{hyperref}
\newcommand{\ie}{i.e.\ }
\newcommand{\cf}{cf.\ }
\newcommand{\resp}{respectively\ }
\newcommand{\Subsection}{\subsection}
\providecommand{\nobreakdash}{\nobreak\hbox{-}}
\setlength{\emergencystretch}{2em}
\multlinegap0pt
\usepackage{bm}
\usepackage{bbm}
\usepackage{dsfont}
\usepackage{xspace}
\usepackage{cancel}
\usepackage{mathtools}
\usepackage{amsthm}
\makeatletter
\@ifundefined{prop}{\newtheorem{prop}{Prop}}{}
\makeatother
\newcommand{\dd}{\mathrm{d}}
\newcommand{\evdensity}{\varrho}
\newcommand{\ii}{\mathrm{i}}
\newcommand{\phaseint}{\Phi}
\newcommand{\tailint}{\Xi}
\title[Suleimanov-Talanov self-focusing and the hierarchy of the fo]{Suleimanov-Talanov self-focusing and the hierarchy of the focusing nonlinear Schr\"odinger equation}
\author{Robert J. Buckingham, Robert M. Jenkins,, Peter D. Miller}
\date{Source: arXiv:2507.01393v1 (CC BY 4.0)}
\begin{document}
\maketitle

\noindent\emph{Source and licence.} Suleimanov-Talanov self-focusing and the hierarchy of the focusing nonlinear Schr\"odinger equation. Version arXiv:2507.01393v1 (CC BY 4.0), distributed under the Creative Commons Attribution 4.0 International licence, as recorded in the arXiv licence metadata for this exact version and shown on the abstract page https://arxiv.org/abs/2507.01393v1. Full rights notice: licenses/arxiv-2507.01393-CC-BY-4.0.txt.\par\smallskip

\noindent\textit{Editorial scope.} Counted unit: the Prop whose source label is \texttt{eq:xpm-recover} in the article (source0.tex lines 692--699, source label \texttt{eq:xpm-recover}), delivered together with the article's own complete original proof of it (source0.tex lines 700--743, one \texttt{proof} environment of 44 lines). No mathematics is written, repaired or re-derived: statement and proof are the article's own text line for line. Required context retained uncounted: eq:ZS-phase-integral (lines 605--608); eq:mu-define (lines 616--620); eq:mu-rewrite (lines 671--675). Every definition, display and label of those lines is retained unchanged. Omitted from this package and not redistributed: 1--604, 609--615, 621--670, 676--691, 744--4479. Nothing inside a retained excerpt is omitted or altered: every retained body is delivered exactly as the article's lines are written, so no omission receipt of any kind accompanies this package. Citations: the retained excerpts carry no citation command at all, so no bibliography entry of the article is transcribed and no reference is redistributed. Reference adaptation: none was needed and none was made; every reference inside the delivered unit resolves inside it, so no unresolved reference or citation is produced. Printed slips kept literal: no printed word, symbol, hypothesis or number of the counted statement or of its proof was altered. Layout and class: the article's own class is \texttt{amsart} with packages including fullpage, mathpazo, mathabx, amssymb, bbm, comment, epstopdf, enumerate, float, graphicx, overpic, mathrsfs, mathtools, multirow; the wrapper is the lane's pinned \texttt{amsart} editorial wrapper at 10pt on A4, which restates the theorem-like environments the retained text needs and adds the article's own macro definitions that the retained text needs (\texttt{\textbackslash dd}, \texttt{\textbackslash evdensity}, \texttt{\textbackslash ii}, \texttt{\textbackslash phaseint}, \texttt{\textbackslash tailint}), with extra wrapper packages (bm, bbm, dsfont, xspace, cancel, mathtools). The delivered numbering is the unit's own. Assets: no retained span contains a figure environment, a graphics-inclusion command, a PDF-page inclusion, a table or a TikZ picture; no image file is copied into this package, the package declares no asset dependency and no image file is redistributed with it. Licence: version arXiv:2507.01393v1 of this article is distributed under the Creative Commons Attribution 4.0 International licence, as recorded in the arXiv licence metadata for this exact version and shown on the abstract page https://arxiv.org/abs/2507.01393v1. A full rights notice is delivered at licenses/arxiv-2507.01393-CC-BY-4.0.txt. Characters: every retained excerpt is pure ASCII, so the delivered engine, XeLaTeX with its default Latin Modern text font, reports no missing character.
\begin{equation}
\phaseint(\ii s):=\pi\int_{s}^{A_\mathrm{max}}\evdensity(s')\,\dd s'= \int_{x_-(s)}^{x_+(s)}\sqrt{A(x)^2-s^2}\,\dd x,\quad 0<s<A_\mathrm{max}.
\label{eq:ZS-phase-integral}
\end{equation}

\begin{multline}
\tailint(\ii s):=(x_+(s)+x_-(s))s + \int_{-\infty}^{x_-(s)}\left(\sqrt{s^2-A(x)^2}-s\right)\,\dd x \\
   - \int_{x_+(s)}^{+\infty}\left(\sqrt{s^2-A(x)^2}-s\right)\,\dd x,\quad 0<s<A_\mathrm{max}.
\label{eq:mu-define}
\end{multline}

\begin{align}
\tailint(\ii s)&=(x_+(s)+x_-(s))s + \int_{X_-}^{x_-(s)}\left(\sqrt{s^2-A(x)^2}-s\right)\,\dd x -\int_{x_+(s)}^{X_+}\left(\sqrt{s^2-A(x)^2}-s\right)\,\dd x\nonumber \\
&=(X_++X_-)s + \int_{X_-}^{x_-(s)}\sqrt{s^2-A(x)^2}\,\dd x -\int_{x_+(s)}^{X_+}\sqrt{s^2-A(x)^2}\,\dd x.
\label{eq:mu-rewrite}
\end{align}

\begin{prop}
Let $A$ be a semicircular Klaus-Shaw potential, and let $\phaseint(\lambda)$ and $\tailint(\lambda)$ be defined by \eqref{eq:ZS-phase-integral} and \eqref{eq:mu-define}, respectively. Then, for $0<s<A_\mathrm{max}$, the inverse functions $x_-(s)<x_+(s)$ of $x\mapsto A(x)$ satisfy
\begin{equation}
x_\pm(s)=\frac{1}{\pi}\int_0^s\frac{\displaystyle\frac{\dd}{\dd m}\tailint(\ii m)\,\dd m}{\sqrt{s^2-m^2}}\mp\frac{1}{\pi}\int_s^{A_\mathrm{max}}\frac{\displaystyle\frac{\dd}{\dd m}\phaseint(\ii m)\,\dd m}{\sqrt{m^2-s^2}}.
\label{eq:xpm-recover}
\end{equation}
\label{prop:Psi-mu-invert}
\end{prop}

\begin{proof}
Directly from the definition of $\phaseint$ in \eqref{eq:ZS-phase-integral}, 
we can compute
\begin{equation}
\frac{\dd}{\dd m}\phaseint(\ii m) = -m\int_{x_-(m)}^{x_+(m)}\frac{\dd x}{\sqrt{A(x)^2-m^2}}.
\end{equation}
Therefore, integrating and changing the integration variable from $m$ to $v=m^2$,
\begin{equation}
\int_s^{A_\mathrm{max}}\frac{\displaystyle\frac{\dd}{\dd m}\phaseint(\ii m)\,\dd m}{\sqrt{m^2-s^2}} = 
-\frac{1}{2}\int_{s^2}^{A_\mathrm{max}^2}\int_{x_-(\sqrt{v})}^{x_+(\sqrt{v})}\frac{\dd x}{\sqrt{A(x)^2-v}}\frac{\dd v}{\sqrt{v-s^2}}.
\end{equation}
Exchanging the order of integration,
\begin{equation}
\int_s^{A_\mathrm{max}}\frac{\displaystyle\frac{\dd}{\dd m}\phaseint(\ii m)\,\dd m}{\sqrt{m^2-s^2}} = -\frac{1}{2}\int_{x_-(s)}^{x_+(s)}\int_{s^2}^{A(x)^2}\frac{\dd v}{\sqrt{A(x)^2-v}\sqrt{v-s^2}}\,\dd x.
\end{equation}
By an affine transformation taking the integration endpoints $v=s^2$ and $v=A(x)^2$ to $\pm 1$, one sees that the inner integral over $v$ is actually independent of $s^2$ and $A(x)^2$, and evaluates to $\pi$. 
Therefore by evaluating the outer integral we obtain
\begin{equation}
\int_s^{A_\mathrm{max}}\frac{\displaystyle\frac{\dd}{\dd m}\phaseint(\ii m)\,\dd m}{\sqrt{m^2-s^2}} = -\frac{\pi}{2}(x_+(s)-x_-(s)).
\label{eq:Psi-integral}
\end{equation}
Similarly, since for $A(x)$ a semicircular Klaus-Shaw potential we can express $\tailint$ using \eqref{eq:mu-rewrite},
we take a derivative to obtain
\begin{align}
\frac{\dd}{\dd m}\tailint(\ii m) &= \frac{\dd}{\dd m}\left[(X_++X_-)m + \int_{X_-}^{x_-(m)}\sqrt{m^2-A(x)^2}\,\dd x -\int_{x_+(m)}^{X_+}\sqrt{m^2-A(x)^2}\,\dd x\right] \nonumber \\
&= X_++X_- + m\int_{X_-}^{x_-(m)}\frac{\dd x}{\sqrt{m^2-A(x)^2}} -m\int_{x_+(m)}^{X_+}\frac{\dd x}{\sqrt{m^2-A(x)^2}}.
\end{align}
So, integrating and changing the integration variable from $m$ to $v=m^2$,
\begin{multline}
\int_0^s\frac{\displaystyle\frac{\dd}{\dd m}\tailint(\ii m)\,\dd m}{\sqrt{s^2-m^2}} = (X_++X_-)\int_0^s\frac{\dd m}{\sqrt{s^2-m^2}} \\
{}+\frac{1}{2}\int_0^{s^2}\int_{X_-}^{x_-(\sqrt{v})}\frac{\dd x}{\sqrt{v-A(x)^2}}\frac{\dd v}{\sqrt{s^2-v}} -\frac{1}{2}\int_0^{s^2}\int_{x_+(\sqrt{v})}^{X_+}\frac{\dd x}{\sqrt{v-A(x)^2}}\frac{\dd v}{\sqrt{s^2-v}}.
\end{multline}
The integral on the first line evaluates to $\pi/2$ independent of $m$, and on the second line we exchange the order of integration in each integral to obtain
\begin{multline}
\int_0^s\frac{\displaystyle\frac{\dd}{\dd m}\tailint(\ii m)\,\dd m}{\sqrt{s^2-m^2}} =\frac{\pi}{2}(X_++X_-)\\ +
\frac{1}{2}\int_{X_-}^{x_-(s)}\int_{A(x)^2}^{s^2}\frac{\dd v}{\sqrt{v-A(x)^2}\sqrt{s^2-v}}\,\dd x -\frac{1}{2}\int_{x_+(s)}^{X_+}\int_{A(x)^2}^{s^2}\frac{\dd v}{\sqrt{v-A(x)^2}\sqrt{s^2-v}}\,\dd x.
\end{multline}
Again the inner integral evaluates to $\pi$ in each case, so we obtain simply
\begin{equation}
\int_0^s\frac{\displaystyle\frac{\dd}{\dd m}\tailint(\ii m)\,\dd m}{\sqrt{s^2-m^2}} =\frac{\pi}{2}(x_+(s)+x_-(s)).
\label{eq:mu-integral}
\end{equation}
Combining \eqref{eq:Psi-integral} and \eqref{eq:mu-integral} yields \eqref{eq:xpm-recover}.
\end{proof}

\end{document}
